\subsection{Strictly subcritical collisions in the circular-loop family}
\label{sec:circular-loop}
A distinct exact construction avoids velocity-one sites on
the marked three-torus. It uses the known circular-loop
Gowdy solutions, locally of Bianchi type $\mathrm{VII}_0$
\cite[Section~3]{rendall-twisted} \cite{grubisic}.
The information comparison and the nonlinear shift of the
nodal placement follow directly from their radius equation.

Use hyperbolic geodesic-polar coordinates at a target point
$U_0$. For integer $k\ge1$ write the two maps as
\[
 U_\pm(t,\theta)=\exp_{U_0}\!\left(
 R(t)[\pm\cos(k\theta)e_P+\sin(k\theta)e_Q]\right).
\]
Here the two basis vectors are evaluated at $U_0$,
and a signed radius makes the maps smooth at zero.
The circular-loop reduction gives
\begin{equation}
 R_{tt}+t^{-1}R_t+k^2\sinh R\cosh R=0.
 \label{eq:circular-ode}
\end{equation}
The spatial and temporal target derivatives are orthogonal.
The metric potential is independent of $\theta$ and
the same on both branches with a common anchor.
The target differential has rank two where
$R R_t\ne0$.

\begin{proposition}
\label{prop:circular-loop}
Fix any areal-time ratio $\mathcal R>1$.
For arbitrarily small prescribed nonzero logarithmic
initial radius velocity, there are times $a,b$ with
$b/a=\mathcal R$ and a circular-loop pair with
identical induced spatial metrics and identical $Q_t$
at both times and with identical $\lambda$ throughout.
Their initial $P_t$ differ in the marked data.
The shared asymptotic speed is strictly between zero
and one and is spatially constant.
If $\nu_0=ka$ is a simple linear Dirichlet frequency
for the ratio $\mathcal R$, the small-amplitude
placement has the expansion
\[
 \nu(\epsilon)=\nu_0-
       \frac{\nu_0}{3}
        \frac{\int_1^{\mathcal R} r\,u(r)^4\,dr}
             {\int_1^{\mathcal R} r\,u(r)^2\,dr}
              \epsilon^2+O(\epsilon^4),
\]
where $u$ has initial derivative one at $r=1$ and
vanishes at both ends.
\end{proposition}

\begin{proof}
In the coordinate $r=t/a$ prescribe
$R(1)=0$, $R'(1)=\epsilon$ and set $\nu=ka$.
The equation is
$-(rR')'=\nu^2 r\sinh R\cosh R$.
The quotient $R/\epsilon$ extends analytically and
evenly in $\epsilon$ near zero and follows the
linear radial Dirichlet equation at zero amplitude.
For a linear eigenfrequency $\nu_0$, the derivative
of the terminal value with respect to $\nu$ is
nonzero. Green's identity expresses its product
with $\mathcal R u'(\mathcal R)$ as
$2\nu_0\int_1^{\mathcal R}r u^2\,dr$.
The implicit function theorem gives the nearby
exact placement as an even analytic function of
$\epsilon$. Substitute
$R=\epsilon u+\epsilon^3 R_3+O(\epsilon^5)$,
$\nu=\nu_0+\epsilon^2\nu_2+O(\epsilon^4)$.
Projecting the cubic equation against $u$
and using $\sinh R\cosh R=R+\frac23 R^3+O(R^5)$
gives the stated $\nu_2$.

At the nodal times both maps return to $U_0$.
Their $Q_t$ equals $e^{-P_0}R_t\sin(k\theta)$
on those slices, and their $P_t$ are
$\pm R_t\cos(k\theta)$.
The equal metric potentials follow from equal
wave-map energy sources and the zero momentum source.

For the limit as $t\downarrow0$, set
$x=\log(a/t)\ge0$. The radius solves
$R_{xx}+\nu^2e^{-2x}\sinh R\cosh R=0$ and
\[
 \frac{d}{dx}\left[\frac12 R_x^2+
                \frac12\nu^2e^{-2x}\sinh^2R\right]
 =-\nu^2e^{-2x}\sinh^2R .
\]
Hence $|R_x|\le|\epsilon|$ and
$|R|\le|\epsilon|x$ for all $x\ge0$.
For $|\epsilon|<1$ the restoring force is integrable
there, so $R_x$ has a limit.
The limit divided by $\epsilon$ tends to
$-J_0(\nu_0)$ as $\epsilon\to0$: this is the
normalized linear Bessel limit, and the force
integral converges under the uniform bound just
obtained and dominated convergence.
One can choose a linear Dirichlet frequency with
$J_0(\nu_0)\ne0$ for any fixed ratio.
Indeed $\nu_j=j\pi/(\mathcal R-1)+O(j^{-1})$,
whereas the positive zeros of $J_0$ have expansion
$(m-\tfrac14)\pi+O(m^{-1})$.
If every sufficiently large $\nu_j$ were such a zero,
successive differences would force
$1/(\mathcal R-1)$ to be an integer, contradicting
the quarter-phase shift. Thus the speed is nonzero
for small enough nonzero amplitude and strictly
below one, with the same value at every angle.
\end{proof}

In this circular-loop example a reflection of the
quotient circle exchanges the two histories.
The collision is stated for the marked record.
The separate two-mode Ernst construction supplies
non-equivalence under boundary-preserving
torus-equivariant isometries when quotient
reflections are also allowed.
